%%%%%%%%%%%%%%%%%%%%%%%%%%%%%%%%%%%%%%%%%%%%%%%%%%%%%%%%%%%%%%%%%%%%%%
%
% Apendice SBSeg'2025 - Filo-Transformer
% Artigo 10657
% Versao 1.0 - Julho 2025
%
%%%%%%%%%%%%%%%%%%%%%%%%%%%%%%%%%%%%%%%%%%%%%%%%%%%%%%%%%%%%%%%%%%%%%%

\documentclass[12pt]{article}

\usepackage{sbc-template}
\usepackage{graphicx,url}
\usepackage[brazil]{babel}   
\usepackage[utf8]{inputenc}
\usepackage{listings}
\usepackage{color}

% Configuracao para codigos
\definecolor{codegreen}{rgb}{0,0.6,0}
\definecolor{codegray}{rgb}{0.5,0.5,0.5}
\definecolor{codepurple}{rgb}{0.58,0,0.82}
\definecolor{backcolour}{rgb}{0.95,0.95,0.92}

\lstdefinestyle{mystyle}{
    backgroundcolor=\color{backcolour},   
    commentstyle=\color{codegreen},
    keywordstyle=\color{magenta},
    numberstyle=\tiny\color{codegray},
    stringstyle=\color{codepurple},
    basicstyle=\ttfamily\footnotesize,
    breakatwhitespace=false,         
    breaklines=true,                 
    captionpos=b,                    
    keepspaces=true,                 
    numbers=left,                    
    numbersep=5pt,                  
    showspaces=false,                
    showstringspaces=false,
    showtabs=false,                  
    tabsize=2,
    literate={á}{{\'a}}1 {ã}{{\~a}}1 {é}{{\'e}}1 {ê}{{\^e}}1 {í}{{\'i}}1 {ó}{{\'o}}1 {õ}{{\~o}}1 {ú}{{\'u}}1 {ç}{{\c{c}}}1
}

\lstset{style=mystyle}

\sloppy

\title{Ap\^endice SBSeg'2025 - Artigo \#10657\\
\large Filo-Transformer: Um modelo baseado em Grafo de Alinhamento de \'Arvores Filogen\'eticas e Transformers para Identifica\c{c}\~ao de Rumores e Fake News}

\author{Acauan C. Ribeiro\inst{1,2}, Eduardo L. Feitosa\inst{1}, Andr\'e Carvalho\inst{1}}

\address{Instituto de Computa\c{c}\~ao -- Universidade Federal do Amazonas (UFAM)\\
  Manaus -- AM -- Brasil
\nextinstitute
  Departamento de Ci\^encia da Computa\c{c}\~ao -- Universidade Federal de Roraima (UFRR)\\
  Boa Vista -- RR -- Brasil
  \email{acauan.ribeiro@ufrr.br, \{efeitosa,andre\}@icomp.ufam.edu.br}
}

\begin{document} 

\maketitle

\begin{resumo} 
Este trabalho prop\~oe o Filo-Transformer, um modelo inovador que integra an\'alise filogen\'etica textual com a arquitetura FT-Transformer para detec\c{c}\~ao de fake news. Atrav\'es da constru\c{c}\~ao de Grafos de Alinhamento de \'Arvores (TAGs) e extra\c{c}\~ao de 70 caracter\'isticas filogen\'eticas avan\c{c}adas das cascatas de propaga\c{c}\~ao, o modelo aprende automaticamente a import\^ancia relativa entre features sem\^anticas e estruturais. Experimentos otimizados no dataset PHEME demonstram melhorias consistentes, com AUC de 0.9237 comparado a 0.9044 do baseline (+2.13\%), validando a import\^ancia das caracter\'isticas filogen\'eticas. Este ap\^endice fornece informa\c{c}\~oes detalhadas para a reprodu\c{c}\~ao completa dos experimentos com configura\c{c}\~ao otimizada.
\end{resumo}

\section{Selos Considerados}

Os selos considerados para este artefato s\~ao: \textbf{Artefatos Dispon\'iveis (SeloD)}, \textbf{Artefatos Funcionais (SeloF)}, \textbf{Artefatos Sustent\'aveis (SeloS)} e \textbf{Experimentos Reprodut\'iveis (SeloR)}.

O artefato atende todos os requisitos atrav\'es de:
\begin{itemize}
    \item \textbf{SeloD}: C\'odigo dispon\'ivel em reposit\'orio GitHub est\'avel com README completo
    \item \textbf{SeloF}: Scripts funcionais com teste m\'inimo e depend\^encias documentadas
    \item \textbf{SeloS}: C\'odigo modularizado, documentado e de f\'acil compreens\~ao
    \item \textbf{SeloR}: Script principal \texttt{main\_experiment.py} com hiperpar\^ametros otimizados
\end{itemize}

\section{Informa\c{c}\~oes B\'asicas}

\subsection{Acesso ao Artefato}

O artefato est\'a dispon\'ivel no GitHub:
\begin{center}
\url{https://github.com/filotransformer/sbseg}
\end{center}

\noindent\textbf{Estrutura do Reposit\'orio:}
\begin{itemize}
    \item \texttt{datasets/} -- Dataset PHEME inclu\'ido (25.5MB compactado)
    \item \texttt{scripts/} -- Scripts principais do modelo e experimentos
    \item \texttt{visualizations/} -- Visualiza\c{c}\~oes geradas
    \item \texttt{results/} -- Resultados dos experimentos
    \item \texttt{README.md} -- Documenta\c{c}\~ao completa
    \item \texttt{SETUP\_WSL.md} -- Guia para Windows/WSL
    \item \texttt{DOCUMENTATION.md} -- Documenta\c{c}\~ao t\'ecnica
\end{itemize}

\subsection{Ambiente de Execu\c{c}\~ao}

\textbf{Hardware M\'inimo:}
\begin{itemize}
    \item CPU: 4 cores (Intel i5 ou equivalente)
    \item RAM: 8GB (16GB recomendado)
    \item Disco: 10GB livres
    \item GPU: Opcional (CUDA 11.0+ para acelera\c{c}\~ao)
\end{itemize}

\textbf{Software Necess\'ario:}
\begin{itemize}
    \item Sistema Operacional: Linux (Ubuntu 20.04+), macOS ou Windows com WSL2
    \item Python: 3.8 ou superior
    \item Git: Para clonar o reposit\'orio
\end{itemize}

\subsection{Depend\^encias}

As principais depend\^encias Python com vers\~oes testadas:

\begin{verbatim}
torch>=2.2.0,<2.8.0
numpy>=1.24.0,<2.0.0
pandas>=2.0.0,<3.0.0
scikit-learn>=1.3.0,<2.0.0
matplotlib>=3.7.0,<4.0.0
seaborn>=0.12.0,<1.0.0
plotly>=5.16.0,<6.0.0
tqdm>=4.66.0,<5.0.0
transformers>=4.30.0,<5.0.0
sentence-transformers>=2.2.0,<3.0.0
\end{verbatim}

\textbf{Dataset PHEME:} O dataset j\'a est\'a inclu\'ido no reposit\'orio em formato compactado (\texttt{phemernrdataset.tar.bz2}, 25.5MB). N\~ao s\~ao necess\'arios downloads externos.

\subsection{Preocupa\c{c}\~oes com Seguran\c{c}a}

Este artefato \textbf{n\~ao apresenta riscos de seguran\c{c}a}:
\begin{itemize}
    \item N\~ao realiza conex\~oes de rede externas
    \item N\~ao modifica arquivos do sistema
    \item Processa apenas dados locais fornecidos
    \item N\~ao executa c\'odigo externo ou comandos do sistema
    \item Todo processamento \'e contido no ambiente Python
\end{itemize}

\section{Instala\c{c}\~ao}

\subsection{Instala\c{c}\~ao Padr\~ao (Linux/macOS)}

\begin{verbatim}
# 1. Clonar repositorio
git clone https://github.com/filotransformer/sbseg.git
cd sbseg

# 2. Criar ambiente virtual
python3 -m venv venv
source venv/bin/activate

# 3. Instalar dependencias
pip install --upgrade pip
pip install -r requirements.txt

# 4. Preparar dataset (apenas primeira vez)
python scripts/prepare_dataset.py
\end{verbatim}

\subsection{Instala\c{c}\~ao no Windows (WSL2)}

Para Windows, \'e necess\'ario usar WSL2:

\begin{verbatim}
# No PowerShell como Administrador
wsl --install -d Ubuntu-24.04
\end{verbatim}

Ap\'os reiniciar, dentro do WSL2:
\begin{verbatim}
# Instalar pacotes do sistema
sudo apt update
sudo apt install python3 python3-pip python3-venv git build-essential

# Seguir passos da instalacao padrao
\end{verbatim}

\textbf{Nota:} Consulte \texttt{SETUP\_WSL.md} para guia detalhado.

\section{Teste M\'inimo}

Execute o teste m\'inimo para verificar a instala\c{c}\~ao:

\begin{verbatim}
python scripts/quick_test.py
\end{verbatim}

\textbf{Sa\'ida esperada:}
\begin{verbatim}
==================================================
TESTE MINIMO - FILO-TRANSFORMER
SBSeg 2025 - Avaliacao de Artefatos
==================================================

[INFO] Verificando instalacao...
OK PyTorch instalado corretamente
OK Transformers instalado corretamente
OK Dataset PHEME encontrado
OK Estrutura de diretorios correta

[INFO] Executando teste minimo...
OK Processamento de amostra: OK
OK Extracao de features: OK
OK Modelo carregado: OK

[INFO] Instalacao verificada com sucesso!
\end{verbatim}

\textbf{Tempo esperado:} $<$ 30 segundos\\
\textbf{Recursos:} $<$ 1GB RAM

\section{Experimentos}

\subsection{Experimento Principal Otimizado}

\textbf{Recomenda\c{c}\~ao:} Execute o experimento principal otimizado que demonstra claramente a superioridade do Filo-Transformer:

\begin{verbatim}
python scripts/main_experiment.py
\end{verbatim}

Este script utiliza:
\begin{itemize}
    \item Hiperpar\^ametros otimizados via busca bayesiana
    \item 70 features filogen\'eticas TAGs (Tree Alignment Graphs)
    \item Valida\c{c}\~ao cruzada 5-fold
    \item Fus\~ao adaptativa via gate mechanism
\end{itemize}

\textbf{Tempo estimado:} 15-20 minutos (2x mais r\'apido com GPU)\\
\textbf{Recursos m\'aximos:} 4GB RAM, 1GB disco

\subsection{Reivindica\c{c}\~ao Principal: Superioridade do Filo-Transformer com TAGs}

\textbf{Descri\c{c}\~ao:} O Filo-Transformer com 70 features TAGs supera o baseline em todas as m\'etricas, demonstrando melhoria de $\sim$2-5\% em AUC.

\textbf{Comando:}
\begin{verbatim}
python scripts/main_experiment.py
\end{verbatim}

\textbf{Configura\c{c}\~ao otimizada:}
\begin{itemize}
    \item \texttt{BATCH\_SIZE}: 16
    \item \texttt{LEARNING\_RATE}: 3e-5
    \item \texttt{D\_MODEL}: 256
    \item \texttt{N\_HEADS}: 8
    \item \texttt{N\_LAYERS}: 3
    \item \texttt{DROPOUT}: 0.2
    \item \texttt{WEIGHT\_DECAY}: 0.01
\end{itemize}

\textbf{Recursos esperados:}
\begin{itemize}
    \item RAM: 8GB
    \item GPU: Opcional (2x mais r\'apido)
    \item Tempo: 15-20 minutos
    \item Disco: 1GB
\end{itemize}

\textbf{Resultado esperado:}
\begin{verbatim}
==================================================================
EXPERIMENTO PRINCIPAL - FILO-TRANSFORMER vs BASELINE
==================================================================

Carregando dataset PHEME com TAGs...
  Total de cascatas: 5802
  Features semanticas: (5802, 384)
  Features filogeneticas TAGs: (5802, 70)

[... Treinamento 5-fold ...]

RESULTADOS FINAIS:

BASELINE (apenas features semanticas):
  AUC: 0.9044 (±0.0076)

FILO-TRANSFORMER (semanticas + filogeneticas TAGs):
  AUC: 0.9237 (±0.0062)

MELHORIA DO FILO-TRANSFORMER:
  AUC: +2.13%

ANALISE DE PESOS DE FUSAO:
  Peso medio semantico: 34.2%
  Peso medio filogenetico: 65.8%
\end{verbatim}

\subsection{Reivindica\c{c}\~ao \#2: Superioridade do Filo-Transformer}

\textbf{Descri\c{c}\~ao:} O Filo-Transformer alcan\c{c}a AUC de 0.9071, superando o baseline de 0.8882 (+1.89\%).

\textbf{Comando para teste r\'apido (5 epochs, 3 folds):}
\begin{verbatim}
python scripts/pheme_real_cascades_experiment.py \
    --seed 42 --folds 3 --epochs 5
\end{verbatim}

\textbf{Comando para reprodu\c{c}\~ao completa do artigo:}
\begin{verbatim}
python scripts/pheme_real_cascades_experiment.py \
    --seed 42 --folds 5 --epochs 30
\end{verbatim}

\textbf{Configura\c{c}\~ao:}
\begin{itemize}
    \item Arquivo: \texttt{scripts/pheme\_real\_cascades\_experiment.py}
    \item Linha 45: \texttt{num\_epochs = 30} (artigo) ou 5 (teste)
    \item Flags: \texttt{--seed 42 --folds 5}
\end{itemize}

\textbf{Recursos esperados:}
\begin{itemize}
    \item RAM: 4GB
    \item Tempo: 2-3 minutos (teste) ou 2-3 horas (completo)
    \item GPU: Opcional (2x mais r\'apido se dispon\'ivel)
\end{itemize}

\textbf{Resultado esperado:}
\begin{verbatim}
=== Resultados Finais (5-Fold CV) ===
Baseline (Semantico):
  AUC: 0.8882 +/- 0.0076
  Accuracy: 0.8671 +/- 0.0089

Filo-Transformer:
  AUC: 0.9071 +/- 0.0069
  Accuracy: 0.8702 +/- 0.0082

Melhoria AUC: +1.89%
\end{verbatim}

\subsection{Reivindica\c{c}\~ao \#3: Aprendizado Autom\'atico de Pesos}

\textbf{Descri\c{c}\~ao:} O modelo aprende automaticamente a priorizar features filogen\'eticas (65\%) sobre sem\^anticas (35\%).

\textbf{Comando:}
\begin{verbatim}
python scripts/pheme_real_cascades_experiment.py \
    --analyze-weights
\end{verbatim}

\textbf{Recursos esperados:}
\begin{itemize}
    \item RAM: 2GB
    \item Tempo: $<$ 1 minuto
\end{itemize}

\textbf{Resultado esperado:}
\begin{verbatim}
=== Analise de Pesos de Fusao ===
Media geral:
- Peso Semantico: 35.0% +/- 1.2%
- Peso Filogenetico: 65.0% +/- 1.2%
\end{verbatim}

\subsection{Reivindica\c{c}\~ao \#4: Valida\c{c}\~ao de Hip\'oteses}

\textbf{Descri\c{c}\~ao:} Valida\c{c}\~ao estat\'istica das 4 hip\'oteses do artigo sobre caracter\'isticas filogen\'eticas.

\textbf{Comando:}
\begin{verbatim}
python scripts/hypothesis_validation_viz.py
\end{verbatim}

\textbf{Recursos esperados:}
\begin{itemize}
    \item RAM: 3GB
    \item Tempo: 1-2 minutos
    \item Disco: 100MB (visualiza\c{c}\~oes)
\end{itemize}

\textbf{Resultado esperado:}
\begin{verbatim}
=== Validacao de Hipoteses ===
H2.1 - Terminal leaves: p < 0.001 OK
H3.2 - Modelos filogeneticos: p < 0.001 OK
H4.2 - Estrutura correlaciona: p < 0.05 OK
H5.2 - Perfis verificados: p < 0.001 OK

Visualizacoes salvas em: visualizations/hypothesis/
\end{verbatim}

\subsection{Reprodu\c{c}\~ao Completa (Opcional)}

Para executar todos os experimentos do artigo:

\begin{verbatim}
bash scripts/reproduce_all.sh
\end{verbatim}

Este script executa automaticamente:
\begin{enumerate}
    \item Processamento com TAGs (se necess\'ario)
    \item Experimento principal otimizado
    \item An\'alise de pesos de fus\~ao
    \item Valida\c{c}\~ao de hip\'oteses
    \item Gera\c{c}\~ao de todas as visualiza\c{c}\~oes
\end{enumerate}

\textbf{Tempo total:} 30-40 minutos

\section{Informa\c{c}\~oes Adicionais}

\subsection{Estrutura Modular do C\'odigo}

O c\'odigo est\'a organizado em m\'odulos independentes:
\begin{itemize}
    \item \texttt{main\_experiment.py}: \textbf{Experimento principal otimizado}
    \item \texttt{hyperparameter\_optimization.py}: Busca autom\'atica de hiperpar\^ametros
    \item \texttt{process\_pheme\_with\_tags.py}: Extra\c{c}\~ao de 70 features TAGs
    \item \texttt{pheme\_real\_cascades\_experiment\_tags.py}: Modelo com TAGs
    \item \texttt{ft\_transformer.py}: Arquitetura base FT-Transformer
    \item \texttt{hypothesis\_validation\_viz.py}: Valida\c{c}\~oes estat\'isticas
\end{itemize}

\subsection{Arquitetura Otimizada do Filo-Transformer}

O modelo utiliza arquitetura otimizada com:
\begin{itemize}
    \item \textbf{Gate Mechanism}: Fus\~ao adaptativa de features
    \item \textbf{Attention Layers}: 8 cabe\c{c}as, 3 camadas
    \item \textbf{Regulariza\c{c}\~ao}: Dropout 0.2, Weight Decay 0.01
    \item \textbf{Otimiza\c{c}\~ao}: AdamW com OneCycleLR
    \item \textbf{Label Smoothing}: 0.1 para robustez
\end{itemize}

\subsection{Caracter\'isticas Filogen\'eticas TAGs (70 features)}

O sistema extrai 70 caracter\'isticas avan\c{c}adas usando Tree Alignment Graphs:

\textbf{Centralidade e Import\^ancia (3):}
\begin{itemize}
    \item PageRank, Closeness, Betweenness centrality
\end{itemize}

\textbf{Estrutura da Cascata (15):}
\begin{itemize}
    \item Profundidade, largura, n\'os por n\'ivel
    \item Folhas terminais, ramifica\c{c}\~ao
\end{itemize}

\textbf{Din\^amica Temporal (8):}
\begin{itemize}
    \item Velocidade de propaga\c{c}\~ao, burst patterns
    \item Tempo de vida, lat\^encia entre n\'iveis
\end{itemize}

\textbf{Caracter\'isticas de Usu\'arios (10):}
\begin{itemize}
    \item Diversidade, taxa de verifica\c{c}\~ao
    \item Influ\^encia agregada, engajamento
\end{itemize}

\textbf{An\'alise de Muta\c{c}\~ao (12):}
\begin{itemize}
    \item Taxa de muta\c{c}\~ao textual
    \item Entropia ancestral, similaridade sem\^antica
\end{itemize}

\textbf{Detec\c{c}\~ao de Comunidades (8):}
\begin{itemize}
    \item Modularidade, grau de recombina\c{c}\~ao
    \item Coeficiente de clustering
\end{itemize}

\textbf{Estat\'isticas Agregadas (14):}
\begin{itemize}
    \item M\'etricas estat\'isticas e distribui\c{c}\~oes
\end{itemize}

\subsection{Troubleshooting}

\textbf{Problema:} Erro de mem\'oria\\
\textbf{Solu\c{c}\~ao:} Reduza batch\_size ou use configura\c{c}\~ao de teste

\textbf{Problema:} Dataset n\~ao encontrado\\
\textbf{Solu\c{c}\~ao:} Execute \texttt{python scripts/prepare\_dataset.py}

\textbf{Problema:} Depend\^encias incompat\'iveis\\
\textbf{Solu\c{c}\~ao:} Use Python 3.8+ e pip atualizado

\subsection{Suporte e Contato}

Para d\'uvidas ou problemas:
\begin{itemize}
    \item Issues: \url{https://github.com/filotransformer/sbseg/issues}
    \item Email: acauan.ribeiro@ufrr.br
\end{itemize}

\subsection{Cita\c{c}\~ao}

Se utilizar este artefato, por favor cite:
\begin{verbatim}
@inproceedings{ribeiro2025filotransformer,
  title={Filo-Transformer: Um modelo baseado em Grafo 
         de Alinhamento de Arvores Filogeneticas e 
         Transformers para Identificacao de Rumores 
         e Fake News},
  author={Ribeiro, Acauan C. and Feitosa, Eduardo L. 
          and Carvalho, Andre},
  booktitle={SBSeg 2025},
  year={2025}
}
\end{verbatim}

\end{document}